\chapter{Modelli Matematici del Machine Learning e Principi Computazionali}
\label{chap:modelli-computazionali}

\FloatBarrier
\section{I Due Pilastri della Matematica Computazionale nel Machine Learning}

Il trattamento computazionale dell'apprendimento automatico e dell'analisi di grandi moli di dati si fonda su due discipline matematiche complementari e profondamente interconnesse~\cite{strang2019linear}:
\begin{itemize}
    \item \textbf{Algebra Lineare Numerica:} fornisce il linguaggio per descrivere le trasformazioni geometriche, proiettare i dati su sottospazi a dimensionalit\`a ridotta, calcolare decomposizioni matriciali fondamentali (SVD, QR, Cholesky) e quantificare la propagazione degli errori floating-point tramite il numero di condizionamento~\cite{trefethen1997numerical,demmel1997applied}.
    \item \textbf{Ottimizzazione Matematica:} formalizza il processo decisionale attraverso la formulazione di problemi vincolati e non vincolati, studia le condizioni analitiche di ottimalit\`a (teoremi di Fermat e condizioni di Karush-Kuhn-Tucker), investiga la dualit\`a lagrangiana e progetta algoritmi iterativi di discesa scalabili~\cite{boyd2004convex,nocedal2006numerical,bazaraa2006nonlinear}.
\end{itemize}

Nelle applicazioni reali di Data Science, nessun modello di ottimizzazione pu\`o prescindere da algoritmi di algebra lineare stabili ed efficienti, poich\'e la risoluzione della direzione di discesa ad ogni passo iterativo si riconduce sistematicamente alla risoluzione di un sistema algebrico lineare.

\begin{figure}[H]
\centering
\resizebox{0.95\textwidth}{!}{%
\begin{tikzpicture}[font=\small]

    % Scheda Sinistra: anchor=east a (-0.3, 0) per garantire un canale di separazione
    \node[
        draw=blue!25,
        fill=blue!1.5,
        rounded corners=8pt,
        text width=7.1cm,
        inner sep=10pt,
        align=left,
        anchor=east
    ] (left) at (-0.3, 0) {
        {\large\bfseries\color{blue!80!black} Algebra Lineare Numerica}\\[2pt]
        {\scriptsize\bfseries\color{blue!60!black}IL MOTORE COMPUTAZIONALE}\\[8pt]
        \footnotesize
        \textbf{Geometria e Rappresentazione:}\\
        Basi ortonormali, sottospazi, proiettori ortogonali.\\[5pt]
        \textbf{Fattorizzazioni Matriciali:}\\
        Valori singolari (SVD), spettrale, QR, Cholesky.\\[5pt]
        \textbf{Stabilit\`a e Trattamento Errori:}\\
        Condizionamento $\kappa(A)$, propagazione all'indietro.\\[5pt]
        \textbf{Algoritmi su Larga Scala:}\\
        Solutori per matrici sparse, metodi di Krylov/Lanczos.
    };

    % Scheda Destra: anchor=west a (+0.3, 0) con separazione netta di 0.6cm da sinistra
    \node[
        draw=red!25,
        fill=red!1.5,
        rounded corners=8pt,
        text width=7.1cm,
        inner sep=10pt,
        align=left,
        anchor=west
    ] (right) at (0.3, 0) {
        {\large\bfseries\color{red!80!black} Ottimizzazione Matematica}\\[2pt]
        {\scriptsize\color{red!60!black}\bfseries IL FORMALISMO DECISIONALE}\\[8pt]
        \footnotesize
        \textbf{Formulazione dell'Apprendimento:}\\
        Rischio empirico, regolarizzazione, vincoli decisionali.\\[5pt]
        \textbf{Condizioni di Ottimalit\`a:}\\
        Punti stazionari, convessit\`a globale, condizioni KKT.\\[5pt]
        \textbf{Teoria della Dualit\`a:}\\
        Duale di Lagrange, gap di dualit\`a, metodi kernel.\\[5pt]
        \textbf{Metodi di Discesa e Convergenza:}\\
        Gradiente, Newton, Quasi-Newton e direzioni scalate.
    };

    % Barra superiore di inquadramento ancorata con anchor=south sopra i pilastri
    \node[
        draw=gray!30,
        fill=gray!4,
        rounded corners=6pt,
        text width=15.3cm,
        inner sep=8pt,
        align=center,
        anchor=south
    ] (top) at ($(left.north)!0.5!(right.north) + (0, 0.7cm)$) {
        {\large\bfseries\color{gray!90!black} FONDAMENTI DI COMPUTATIONAL MATHEMATICS FOR DATA ANALYSIS}\\[2pt]
        {\footnotesize\color{gray!70!black}I due fondamenti complementari dell'apprendimento automatico moderno}
    };

    % Pannello inferiore ancorato con anchor=north sotto i pilastri
    \node[
        draw=teal!30,
        fill=teal!2,
        rounded corners=6pt,
        text width=15.3cm,
        inner sep=9pt,
        align=center,
        anchor=north
    ] (bot) at ($(left.south)!0.5!(right.south) + (0, -0.7cm)$) {
        {\bfseries\color{teal!80!black} Punto di Confluenza Operativo: Il Ciclo Algoritmico del Machine Learning}\\[4pt]
        \parbox{14.8cm}{\footnotesize\centering
        L'\textbf{ottimizzazione} formula il problema (\textit{cosa} minimizzare e sotto quali vincoli) e genera a ogni passo un sistema direzionale $\nabla^2 f(x) \Delta x = -\nabla f(x)$.\\
        L'\textbf{algebra lineare numerica} fornisce gli algoritmi efficienti e stabili per risolverlo su miliardi di parametri (\textit{come} calcolarlo computazionalmente).}
    };

    % Linee di accento laterali sui pilastri
    \draw[line width=4pt, color=blue!70!black, line cap=round] ($(left.north west)+(2pt,-5pt)$) -- ($(left.south west)+(2pt,5pt)$);
    \draw[line width=4pt, color=red!70!black, line cap=round] ($(right.north west)+(2pt,-5pt)$) -- ($(right.south west)+(2pt,5pt)$);

\end{tikzpicture}%
}
\caption{Quadro concettuale sinergico: integrazione complementare tra Algebra Lineare Numerica e Ottimizzazione Matematica nel Machine Learning.}
\label{fig:architettura-pilastri}
\end{figure}

\FloatBarrier
\section{La Filosofia della Modellazione Matematica: Brunelleschi e Galileo}

\subsection{La Nascita del Modello: Dalla Prospettiva al Calcolo}
Il processo di estrazione di conoscenza da grandi volumi di dati risponde a un'esigenza epistemologica precisa: i dati grezzi sono disorganizzati, cognitivamente inaccessibili all'uomo e incapaci di prendere decisioni autonome.

\begin{definition}[Modello Computazionale]
L'\textbf{apprendimento (machine learning)} consiste nell'estrarre regolarit\`a strutturali da insiemi di osservazioni empiriche per sintetizzarle in un artefatto matematico astratto (\textbf{modello}), manipulabile computazionalmente e atto a generare predizioni o decisioni operative.
\end{definition}

Questa impostazione affonda le sue radici nella storia della scienza:
\begin{itemize}
    \item \textbf{Filippo Brunelleschi e la Prospettiva:} per innalzare l'ardita cupola di Santa Maria del Fiore a Firenze senza armature lignee da terra, Brunelleschi non procedette per tentativi ed errori in cantiere; costru\`i modelli matematici rigorosi (proiezioni geometriche basate sul cono ottico) capaci di descrivere le forze di spinta radiali.
    \item \textbf{L'Economicit\`a della Matematica:} costruire modelli fisici in scala (es. gallerie del vento in aerodinamica) comporta costi materiali ingenti. \textbf{La matematica costa infinitamente meno del legno o dell'acciaio}: una formula o un algoritmo consentono di simulare miliardi di scenari operativi a costo computazionale trascurabile.
    \item \textbf{Galileo Galilei e l'Astrazione:} la fisica moderna nasce nel momento in cui si smette di descrivere passivamente la realt\`a in ogni dettaglio accidentale e si formula una legge astratta parsimoniosa espressa in linguaggio algebrico.
\end{itemize}

\begin{noteblock}{Definizione di Modello Matematico Computazionale}
Un modello non \`e una copia esatta della realt\`a fisica. \`E una \textbf{formalizzazione matematica astratta} volta a catturare le relazioni causali o correlative essenziali del fenomeno, trascurando il rumore al fine di rendere il problema trattabile mediante algoritmi al calcolatore.
\end{noteblock}

\FloatBarrier
\section{Caratterizzazione e Tassonomia dei Modelli}

\subsection{Il Triangolo dei Tre Requisiti Contraddittori}
Nel concepire un modello computazionale si persegue idealmente la massimizzazione congiunta di tre propriet\`a:
\begin{enumerate}
    \item \textbf{Accuratezza (Accurate):} capacit\`a del modello di riprodurre fedelmente la fisica o il comportamento effettivo del sistema analizzato.
    \item \textbf{Economicit\`a Computazionale (Computationally Inexpensive):} semplicit\`a algoritmica, traducibile in basso consumo di memoria e tempo di risoluzione contenuto (complessit\`a lineare $\BigO{n}$ o quasi-lineare).
    \item \textbf{Generalit\`a (General):} attitudine del modello a essere impiegato efficacemente su un ampio spettro di scenari e domini applicativi differenti.
\end{enumerate}

\begin{property}[Principio di Incompatibilit\`a dei Modelli]
\`E teoricamente e praticamente impossibile massimizzare simultaneamente tutti e tre i requisiti: modelli estremamente fedeli richiedono un numero sterminato di parametri che ne distrugge l'economicit\`a computazionale; modelli parsimoniosi e rapidi sacrificano inevitabilmente l'accuratezza fine o la generalit\`a. La formulazione di un modello \`e sempre un \textbf{compromesso ingegneristico (\textit{trade-off})}.
\end{property}

\begin{figure}[H]
\centering
\resizebox{0.82\textwidth}{!}{%
\begin{tikzpicture}[>=Stealth, scale=1.0]
    % Vertici del triangolo
    \coordinate (A) at (0, 3.8);
    \coordinate (B) at (-4.2, 0);
    \coordinate (C) at (4.2, 0);

    % Lati del triangolo
    \draw[thick, draw=blue!70!black] (A) -- (B) -- (C) -- cycle;
    \fill[blue!4] (A) -- (B) -- (C) -- cycle;

    % Nodi ai vertici
    \node[above=4pt, font=\bfseries\normalsize, text=blue!90!black, align=center] at (A) {Accuratezza\\\scriptsize (Fedelt\`a analitica)};
    \node[below left=2pt, font=\bfseries\normalsize, text=green!60!black, align=center] at (B) {Economicit\`a Computazionale\\\scriptsize (Complessit\`a $\BigO{n}$)};
    \node[below right=2pt, font=\bfseries\normalsize, text=orange!80!black, align=center] at (C) {Generalit\`a\\\scriptsize (Trasferibilit\`a di dominio)};

    % Punto di compromesso interno
    \node[circle, fill=red!80!black, inner sep=2.5pt] (P) at (0, 1.4) {};
    \node[below=5pt, font=\footnotesize\bfseries, text=red!80!black, align=center] at (P) {Compromesso Ingegneristico\\(Punto di equilibrio operativo)};

    % Frecce bidirezionali sui lati
    \draw[<->, dashed, gray!60, thick] (-1.8, 2.0) -- (-0.4, 2.8);
    \draw[<->, dashed, gray!60, thick] (1.8, 2.0) -- (0.4, 2.8);
    \draw[<->, dashed, gray!60, thick] (-2.2, 0.3) -- (2.2, 0.3);
\end{tikzpicture}%
}
\caption{Il triangolo dei requisiti contrastanti nella modellazione computazionale: accuratezza, economicit\`a e generalit\`a.}
\label{fig:triangolo-modelli}
\end{figure}

\subsection{Modelli Analitici vs Modelli Guidati dai Dati (Data-Driven)}
La tabella~\ref{tab:confronto-modelli} riassume la demarcazione strutturale tra le due polarit\`a della modellazione computazionale.

\begin{table}[H]
\centering
\small
\begin{tabularx}{\textwidth}{@{} l Y Y @{}}
\toprule
\textbf{Propriet\`a} & \textbf{Approccio Analitico (White-Box)} & \textbf{Approccio Guidato dai Dati (Black-Box)} \\
\midrule
\textbf{Principio Base} & Modella analiticamente ogni singolo componente fisico e le sue interazioni differenziali. & Non descrive la meccanica interna del sistema; modella direttamente la relazione ingresso-uscita. \\
\addlinespace
\textbf{Struttura Matematica} & Flessibile, non lineare, basata su principi primi fisici e leggi di conservazione. & Standardizzata, rigida, parametrica (es. reti neurali profonde, SVM, iperpiani lineari). \\
\addlinespace
\textbf{Trattamento Parametri} & Pochi parametri calibrabili, dotati di preciso significato fisico intrinseco. & Enorme numero di parametri (da milioni a miliardi), stimati automaticamente via ottimizzazione. \\
\addlinespace
\textbf{Costo di Sviluppo} & Elevato: richiede conoscenza di dominio avanzata e accesso interno al sistema. & Basso/scalabile: necessita di osservazioni storiche e algoritmi di ottimizzazione standardizzati. \\
\bottomrule
\end{tabularx}
\caption{Confronto concettuale tra modellazione analitica basata su principi primi e modellazione empirica guidata dai dati.}
\label{tab:confronto-modelli}
\end{table}

\subsection{Fitting vs Machine Learning: Rischio Empirico e Generalizzazione}
Un modello \`e sempre un'approssimazione (\textit{«the map is not the territory»}):
\begin{itemize}
    \item \textbf{Fitting (Minimizzazione del Rischio Empirico):} ricerca i parametri ottimali che minimizzano lo scostamento esclusivamente sui punti noti del campione di addestramento (\textit{training set}):
    \begin{equation}
        \min_{\mathbf{w}} \mathcal{R}_{\text{emp}}(\mathbf{w}) = \frac{1}{m} \sum_{i=1}^m \ell(f(\mathbf{x}^i; \mathbf{w}), y^i).
    \end{equation}
    \item \textbf{Generalizzazione (Apprendimento):} richiede che $f(\mathbf{x}; \mathbf{w})$ mantenga un errore contenuto su dati non visti durante l'addestramento (\textit{test set}). Un modello iper-parametrizzato che minimizza a zero l'errore empirico incorre nell'\textbf{overfitting}, catturando il rumore anzich\'e il segnale.
\end{itemize}

\FloatBarrier
\section{I Quattro Modelli Matematici Cardine del Corso}

\subsection{Esempio 1: Stima Lineare e Minimi Quadrati}
Dato un insieme di $m$ osservazioni $(\mathbf{x}^i, y^i) \in \R^n \times \R$, si ipotizza una relazione lineare tra ingressi e uscita:
\begin{equation}
    f(\mathbf{x}) = \mathbf{w}^T \mathbf{x} + w_0 = \sum_{j=1}^n w_j x_j + w_0.
\end{equation}
Inglobando il termine di bias $w_0$ nella matrice di progetto estesa $X$:
\begin{equation}
    \mathbf{w}_{\text{ext}} = \begin{bmatrix} w_0 \\ \mathbf{w} \end{bmatrix} \in \R^{n+1}, \quad
    X = \begin{bmatrix} 1 & (\mathbf{x}^1)^T \\ \vdots & \vdots \\ 1 & (\mathbf{x}^m)^T \end{bmatrix} \in \R^{m \times (n+1)}, \quad
    \mathbf{y} = \begin{bmatrix} y^1 \\ \vdots \\ y^m \end{bmatrix} \in \R^m.
\end{equation}
Il problema dei minimi quadrati lineari consiste nel risolvere:
\begin{equation}
    \min_{\mathbf{w}} \mathcal{L}(\mathbf{w}) = \|\mathbf{y} - X\mathbf{w}\|_2^2 = \mathbf{y}^T \mathbf{y} - 2\mathbf{w}^T X^T \mathbf{y} + \mathbf{w}^T X^T X \mathbf{w}.
\end{equation}
Annullando il gradiente rispetto a $\mathbf{w}$ si perviene alla condizione di stazionariet\`a, nota come \textbf{Equazioni Normali di Gauss}:
\begin{equation}
    \nabla \mathcal{L}(\mathbf{w}) = -2 X^T \mathbf{y} + 2 X^T X \mathbf{w} = \mathbf{0} \implies X^T X \mathbf{w} = X^T \mathbf{y}.
\end{equation}
Se la matrice $X$ ha rango pieno per colonne ($\rank(X) = n+1$), il gramiano $X^T X$ \`e simmetrico, strettamente definito positivo e invertibile:
\begin{equation}
    \mathbf{w}^* = (X^T X)^{-1} X^T \mathbf{y}.
\end{equation}

\begin{alertblock}{Regola Aurea Numerica: Mai Calcolare $(X^T X)^{-1}$}
\label{rem:regola-aurea-normal-equations}
Nel calcolo scientifico professionale \textbf{non si forma mai la matrice $X^T X$ n\'e si calcola la sua inversa esplicita}. L'operazione di prodotto moltiplica al quadrato il numero di condizionamento ($\kappa(X^T X) = \kappa(X)^2$), amplificando catastroficamente gli errori di arrotondamento floating-point. La soluzione si calcola stabilmente mediante fattorizzazione QR o decomposizione SVD (operatore backslash in GNU Octave: \texttt{w = X \textbackslash\ y}).
\end{alertblock}

\subsection{Esempio 2: Approssimazione di Basso Rango e Sistemi di Raccomandazione}
Sia data una matrice dati $M \in \R^{n \times m}$ parzialmente osservata (es. $n$ utenti e $m$ film su una piattaforma, dove $M_{ij}$ rappresenta il punteggio assegnato). Si cerca una fattorizzazione di rango ridotto $k \ll \min(n, m)$ tale che $M \approx AB$:
\begin{equation}
    A \in \R^{n \times k}, \quad B \in \R^{k \times m} \implies \min_{A, B} \|M - AB\|_F^2 = \sum_{i=1}^n \sum_{j=1}^m (M_{ij} - (AB)_{ij})^2.
\end{equation}

\begin{figure}[H]
\centering
\resizebox{0.92\textwidth}{!}{%
\begin{tikzpicture}[>=Stealth]
    % Matrice M
    \node[draw=blue!80!black, fill=blue!10, thick, minimum width=3.4cm, minimum height=4.2cm] (M) at (0, 0) {};
    \node[font=\bfseries\Large] at (0, 0) {$M$};
    \node[above=3pt, font=\small] at (M.north) {$m$ colonne};
    \node[left=3pt, font=\small] at (M.west) {$n$ righe};

    \node[font=\Large] at (2.4, 0) {$\approx$};

    % Matrice A
    \node[draw=red!80!black, fill=red!10, thick, minimum width=1.1cm, minimum height=4.2cm] (A) at (3.8, 0) {};
    \node[font=\bfseries\large] at (3.8, 0) {$A$};
    \node[above=3pt, font=\small] at (A.north) {$k$};
    \node[below=3pt, font=\footnotesize, align=center] at (A.south) {Profili riga\\($n \times k$)};

    \node[font=\large] at (4.8, 0) {$\times$};

    % Matrice B
    \node[draw=green!70!black, fill=green!10, thick, minimum width=3.4cm, minimum height=1.1cm] (B) at (7.0, 0) {};
    \node[font=\bfseries\large] at (7.0, 0) {$B$};
    \node[right=3pt, font=\small] at (B.east) {$k$};
    \node[above=3pt, font=\small] at (B.north) {$m$ colonne};
    \node[below=3pt, font=\footnotesize, align=center] at (B.south) {Profili colonna\\($k \times m$)};
\end{tikzpicture}%
}
\caption{Fattorizzazione matriciale a basso rango: compressione da $nm$ locazioni a $k(n + m)$.}
\label{fig:low-rank-matrix}
\end{figure}

\subsubsection{Compressione di Memoria e Propriet\`a Spettrali}
Memorizzare $M$ densamente richiede $n \cdot m$ locazioni; memorizzare $A$ e $B$ richiede $k(n + m)$ valori. Ad esempio, per $n = 10^6$, $m = 10^5$ e rango $k = 10$:
\begin{align}
    \text{Memoria originaria} &= 10^6 \times 10^5 = 10^{11} \text{ celle}, \\
    \text{Fattorizzata} &= 10 \times (10^6 + 10^5) \approx 1.1 \times 10^7 \text{ celle},
\end{align}
ottenendo un \textbf{abbattimento dello spazio occupato del 99.99\%}.

In virt\`u del Teorema di Eckart-Young-Mirsky, la soluzione ottima in norma di Frobenius \`e data dalla Decomposizione a Valori Singolari (SVD) troncata.

\begin{alertblock}{Avvertenza Computazionale su Matrici Sparse}
Se $M$ \`e sparsa, \textbf{\`e tassativo non formare mai $M^T M$ n\'e $M M^T$}. Il prodotto matriciale distrugge la sparsit\`a generando una matrice densa che esaurisce all'istante la RAM di sistema (\textit{Out Of Memory}). Si impiegano esclusivamente metodi iterativi a blocchi basati su prodotti matrice-vettore (metodi di Lanczos).
\end{alertblock}

\subsection{Esempio 3: Support Vector Machines (SVM) --- Margine, Dualit\`a e Kernel}
Nel contesto della classificazione binaria supervisionata con etichette $y^h \in \{-1, +1\}$, un classificatore lineare assegna la classe mediante la regola di decisione $\hat{y} = \operatorname{sign}(\mathbf{w}_+^T \mathbf{x} - w_0)$.

\begin{figure}[H]
\centering
\resizebox{0.82\textwidth}{!}{%
\begin{tikzpicture}[>=Stealth, scale=1.05]
    % Assi coordinati
    \draw[->, thick, gray!60] (-0.5, 0) -- (7.5, 0) node[right, black] {$x_1$};
    \draw[->, thick, gray!60] (0, -0.5) -- (0, 6.6) node[above, black] {$x_2$};
    \draw[help lines, color=gray!15, dashed] (0,0) grid (7.0,6.2);

    % Fascia del margine (sfondo evidenziato)
    \fill[teal!6, opacity=0.7] (0.2, 2.4) -- (3.0, 5.2) -- (5.8, 5.6) -- (1.4, 1.2) -- cycle;

    % Iperpiano di supporto superiore (+1): x2 = x1 + 2.2
    \draw[dashed, thick, draw=blue!75!black] (0.2, 2.4) -- (3.0, 5.2) 
        node[above left=3pt, font=\footnotesize\bfseries, blue!85!black, fill=white, inner sep=1.5pt] {$\mathbf{w}_+^T \mathbf{x} - w_0 = +1$};

    % Iperpiano di separazione ottimo (w^T x - w_0 = 0): x2 = x1 + 1.0
    \draw[very thick, draw=black] (0.8, 1.8) -- (4.8, 5.8) 
        node[above right=3pt, font=\footnotesize\bfseries, black, fill=white, inner sep=1.5pt] {$\mathbf{w}_+^T \mathbf{x} - w_0 = 0$};

    % Iperpiano di supporto inferiore (-1): x2 = x1 - 0.2
    \draw[dashed, thick, draw=red!75!black] (1.4, 1.2) -- (5.8, 5.6) 
        node[right=4pt, font=\footnotesize\bfseries, red!85!black, fill=white, inner sep=1.5pt] {$\mathbf{w}_+^T \mathbf{x} - w_0 = -1$};

    % Punti interni Classe +1 (cerchi blu)
    \foreach \coord in {(0.8, 4.6), (1.4, 5.4), (2.0, 5.8), (0.5, 3.2)} {
        \fill[blue!75!black] \coord circle (3.0pt);
    }

    % Punti interni Classe -1 (quadrati rossi)
    \foreach \coord in {(3.6, 1.6), (4.6, 2.0), (5.4, 2.4), (4.2, 0.8), (5.6, 1.4)} {
        \fill[red!75!black] \coord +(-3.0pt,-3.0pt) rectangle +(3.0pt,3.0pt);
    }

    % Support Vectors (esattamente sugli iperpiani di supporto)
    % Su H+1: (1.6, 3.8) e (2.6, 4.8)
    \fill[blue!75!black] (1.6, 3.8) circle (3.0pt);
    \draw[line width=1.5pt, green!60!black] (1.6, 3.8) circle (6.5pt);

    \fill[blue!75!black] (2.6, 4.8) circle (3.0pt);
    \draw[line width=1.5pt, green!60!black] (2.6, 4.8) circle (6.5pt);

    % Su H-1: (2.8, 2.6) e (3.8, 3.6)
    \fill[red!75!black] (2.8, 2.6) +(-3.0pt,-3.0pt) rectangle +(3.0pt,3.0pt);
    \draw[line width=1.5pt, green!60!black] (2.8, 2.6) circle (6.5pt);

    \fill[red!75!black] (3.8, 3.6) +(-3.0pt,-3.0pt) rectangle +(3.0pt,3.0pt);
    \draw[line width=1.5pt, green!60!black] (3.8, 3.6) circle (6.5pt);

    % Etichette Support Vector con puntatori orizzontali senza incroci
    \node[anchor=east, font=\scriptsize\bfseries, green!50!black, fill=white, inner sep=1.5pt] (sv1) at (0.6, 3.8) {Support Vector};
    \draw[->, thick, green!60!black] (sv1.east) -- (1.35, 3.8);

    \node[anchor=west, font=\scriptsize\bfseries, green!50!black, fill=white, inner sep=1.5pt] (sv2) at (4.8, 3.6) {Support Vector};
    \draw[->, thick, green!60!black] (sv2.west) -- (4.05, 3.6);

    % Vettore normale w_+ ortogonale all'iperpiano di separazione
    \draw[->, ultra thick, purple] (4.4, 5.4) -- (3.7, 6.1);
    \node[above=2pt, font=\small\bfseries, purple, fill=white, inner sep=1pt] at (3.7, 6.1) {$\mathbf{w}_+$};
    \draw[purple!80!black, thin] (4.4, 5.4) +(-0.1, 0.1) -- +(-0.0, 0.2) -- +(0.1, 0.1);

    % Ampiezza del margine (distanza ortogonale tra H-1 e H+1 posizionata in area sgombra)
    \draw[<->, line width=1.3pt, purple!85!black] (2.0, 1.8) -- (0.8, 3.0);
    \node[anchor=north east, font=\footnotesize\bfseries, purple!85!black, fill=white, inner sep=2pt] at (1.2, 2.2) {Margine $\delta = \frac{2}{\|\mathbf{w}_+\|_2}$};

    % Legenda
    \node[draw=gray!30, fill=white, rounded corners=2pt, font=\scriptsize, anchor=south east] at (7.3, 0.2) {
        \begin{tabular}{cl}
        \tikz\fill[blue!75!black] (0,0) circle (2.5pt); & Classe $y = +1$ \\
        \tikz\fill[red!75!black] (-2.5pt,-2.5pt) rectangle (2.5pt,2.5pt); & Classe $y = -1$ \\
        \tikz\draw[line width=1.3pt, green!60!black] (0,0) circle (4pt); & Support Vector
        \end{tabular}
    };
\end{tikzpicture}%
}
\caption{Geometria del margine nelle Support Vector Machines lineari: massimizzazione della distanza ortogonale tra gli iperpiani di supporto.}
\label{fig:svm-margin}
\end{figure}

Massimizzare il margine geometrico $2/\|\mathbf{w}_+\|_2$ equivale a minimizzare il semiquadrato della norma del vettore dei pesi:
\begin{equation}
    \min_{\mathbf{w}_+, w_0} \frac{1}{2} \|\mathbf{w}_+\|_2^2 \qquad \text{s.t.} \quad y^h (\mathbf{w}_+^T \mathbf{x}^h - w_0) \ge 1 \quad \forall h = 1, \dots, m.
\end{equation}

Per dati non perfettamente separabili, si introducono variabili di rilassamento $\xi_h \ge 0$ (\textit{slack variables}), pervenendo al modello a margine morbido (\textbf{Soft Margin}):
\begin{equation}
\begin{aligned}
    \min_{\mathbf{w}_+, w_0, \boldsymbol{\xi}} \quad & \frac{1}{2} \|\mathbf{w}_+\|_2^2 + C \sum_{h=1}^m \xi_h \\
    \text{s.t.} \quad & y^h (\mathbf{w}_+^T \mathbf{x}^h - w_0) \ge 1 - \xi_h, \quad \xi_h \ge 0 \quad \forall h=1, \dots, m.
\end{aligned}
\end{equation}
Il parametro $C > 0$ governa il trade-off bias-varianza: per $C \to \infty$ si impone la separazione rigida, mentre valori finiti consentono violazioni controllate a beneficio della generalizzazione su dati di test.

\subsubsection{Lifting Non Lineare, Dualit\`a e Kernel Trick}
Qualora i punti non siano separabili linearmente nello spazio originario $\R^n$, si applica un'immersione non lineare $\boldsymbol{\phi}: \R^n \to \mathcal{H}$ verso uno spazio delle caratteristiche a dimensionalit\`a superiore (\textit{lifting}):
\begin{equation}
    \mathbf{x} \mapsto \boldsymbol{\phi}(\mathbf{x}).
\end{equation}
Formulando il problema duale di Lagrange (dualit\`a di Wolfe):
\begin{equation}
\begin{aligned}
    \max_{\boldsymbol{\alpha}} \quad & \sum_{h=1}^m \alpha_h - \frac{1}{2} \sum_{h=1}^m \sum_{k=1}^m \alpha_h \alpha_k y^h y^k \langle \boldsymbol{\phi}(\mathbf{x}^h), \boldsymbol{\phi}(\mathbf{x}^k) \rangle_{\mathcal{H}} \\
    \text{s.t.} \quad & \sum_{h=1}^m y^h \alpha_h = 0, \quad 0 \le \alpha_h \le C \quad \forall h=1, \dots, m.
\end{aligned}
\end{equation}
Poich\'e i dati intervengono unicamente tramite prodotti scalari, in virt\`u del Teorema di Mercer si adotta una \textbf{funzione kernel} $K(\mathbf{x}^h, \mathbf{x}^k) = \langle \boldsymbol{\phi}(\mathbf{x}^h), \boldsymbol{\phi}(\mathbf{x}^k) \rangle_{\mathcal{H}}$. Ad esempio, il \textbf{Kernel Gaussiano RBF}:
\begin{equation}
    K(\mathbf{x}, \mathbf{z}) = \exp\left(-\gamma \|\mathbf{x} - \mathbf{z}\|_2^2\right)
\end{equation}
corrisponde a uno spazio di Hilbert $\mathcal{H}$ a dimensione infinita, computabile con costo $\BigO{n}$ senza mai materializzare esplicitamente $\boldsymbol{\phi}(\mathbf{x})$.

\subsection{Esempio 4: Clustering e Algoritmo K-Means}
Dato un dataset non etichettato $X = \{\mathbf{x}^1, \dots, \mathbf{x}^m\} \subset \R^h$, il clustering mira a partizionare $X$ in $k$ cluster disgiunti $\{X_p\}_{p=1}^k$ massimizzando la coesione interna.

\subsubsection{Formulazione Ottimizzatoria Biconvessa}
Introducendo $k$ centroidi $\mathbf{c}^1, \dots, \mathbf{c}^k \in \R^h$ e le variabili binarie di appartenenza $z_{ip} \in \{0, 1\}$:
\begin{equation}
    \min_{\mathbf{c}, Z} \sum_{i=1}^m \sum_{p=1}^k z_{ip} \|\mathbf{c}^p - \mathbf{x}^i\|_2^2 \qquad \text{s.t.} \quad \sum_{p=1}^k z_{ip} = 1 \; \forall i, \quad z_{ip} \in \{0, 1\} \; \forall i,p.
\end{equation}
Il problema \`e fortemente non convesso a causa dell'interezza $z_{ip} \in \{0, 1\}$ e del termine bilineare $z_{ip}\mathbf{c}^p$. Tuttavia, possiede una struttura biconvessa separabile che abilita l'ottimizzazione a \textbf{blocchi alternati (Block Gauss-Seidel)}:
\begin{enumerate}
    \item \textbf{Fissato $\mathbf{c}$, si ottimizza $Z$:} ciascun punto viene assegnato al centroide pi\`u vicino in norma euclidea ($z_{i\bar{p}} = 1 \iff \bar{p} = \arg\min_p \|\mathbf{c}^p - \mathbf{x}^i\|_2^2$).
    \item \textbf{Fissato $Z$, si ottimizza $\mathbf{c}$:} per ciascun cluster $p$, annullando il gradiente quadratico convesso rispetto a $\mathbf{c}^p$:
    \begin{equation}
        2 \sum_{i \in I(p)} (\mathbf{c}^p - \mathbf{x}^i) = \mathbf{0} \implies (\mathbf{c}^p)^* = \frac{1}{|I(p)|} \sum_{i \in I(p)} \mathbf{x}^i.
    \end{equation}
    Il centroide ottimo \`e analiticamente il \textbf{baricentro (media vettoriale)} dei punti appartenenti al cluster.
\end{enumerate}

\begin{algorithm}[H]
\caption{Algoritmo K-Means per Clustering Euclideo}\label{alg:kmeans}
\KwInput{Dataset $X = \{\mathbf{x}^1, \dots, \mathbf{x}^m\} \subset \R^h$, numero di cluster $k$, centroidi iniziali $\mathbf{c}^1, \dots, \mathbf{c}^k$, tolleranza $\epsilon > 0$}
\KwOutput{Centroidi finali $\mathbf{c}^1, \dots, \mathbf{c}^k$ e partizione $\{I(p)\}_{p=1}^k$}
$v \leftarrow \infty$\;
\Repeat{$(v - v_{\mathrm{new}}) \le \epsilon$}{
    \tcp{Fase 1: Assegnamento dei punti ai centroidi}
    \For{$p \leftarrow 1$ \KwTo $k$}{
        $I(p) \leftarrow \emptyset$\;
    }
    \For{$i \leftarrow 1$ \KwTo $m$}{
        $p^* \leftarrow \arg\min_{p \in \{1, \dots, k\}} \|\mathbf{c}^p - \mathbf{x}^i\|_2^2$\;
        $I(p^*) \leftarrow I(p^*) \cup \{i\}$\;
    }
    \tcp{Fase 2: Aggiornamento dei centroidi (baricentro)}
    \For{$p \leftarrow 1$ \KwTo $k$}{
        \If{$I(p) \neq \emptyset$}{
            $\mathbf{c}^p \leftarrow \frac{1}{|I(p)|} \sum_{i \in I(p)} \mathbf{x}^i$\;
        }
    }
    \tcp{Fase 3: Valutazione obiettivo}
    $v_{\mathrm{new}} \leftarrow \sum_{p=1}^k \sum_{i \in I(p)} \|\mathbf{c}^p - \mathbf{x}^i\|_2^2$\;
    $v \leftarrow v_{\mathrm{new}}$\;
}
\Return{$\mathbf{c}^1, \dots, \mathbf{c}^k$}\;
\end{algorithm}

\begin{noteblock}{Perch\'e il Calcolo Scientifico Viola il Paradigma Agile}
Nello sviluppo software tradizionale (Agile), un sistema complesso viene suddiviso in componenti modulari indipendenti modificabili isolatamente. Nel calcolo scientifico ci\`o non \`e possibile:
\begin{itemize}
    \item Se si sostituisce la metrica euclidea $\ell_2$ con la distanza Manhattan $\ell_1$, la forma chiusa del baricentro decade: l'ottimo per $\ell_1$ non \`e pi\`u la media aritmetica, bens\`i la \textbf{mediana componente per componente}.
    \item L'interdipendenza tra la funzione di costo e l'algoritmo di aggiornamento impone una progettazione matematica globale ed integrata.
\end{itemize}
\end{noteblock}

\FloatBarrier
\section{Principi Metodologici dell'Apprendimento Computazionale}

\subsection{La Metafora della ``Magic Academy''}
Per chiarire la distinzione tra l'impiego superficiale di strumenti predefiniti e la reale competenza scientifica, si consideri la metafora della \textbf{Magic Academy}:
\begin{itemize}
    \item \textbf{L'Apprendista Stregone (\textit{Black-Box Practitioner}):} invoca funzioni di alto livello (\texttt{import sklearn}, \texttt{fit()}, \texttt{predict()}) come formule magiche precostituite. Fin quando i dati sono ideali ottiene risultati apparentemente corretti; tuttavia, di fronte a instabilit\`a numerica, perdita di convergenza o mal condizionamento, non possiede gli strumenti per comprendere la causa del fallimento.
    \item \textbf{Il Maestro Mago (\textit{Computational Scientist}):} conosce le propriet\`a spettrali delle matrici e la teoria della convergenza, riconosce quando un problema \`e mal posto, valuta la propagazione degli errori floating-point e sostituisce algoritmi instabili con routine numericamente solide.
\end{itemize}

\subsection{I Tre Principi Guida del Machine Learning}
\begin{enumerate}
    \item \textbf{L'Apprendimento \`e Calcolo Numerico:} Ogni algoritmo di intelligenza artificiale si riduce, in ultima istanza, alla risoluzione di un sistema algebrico lineare o alla minimizzazione di una funzione di costo.
    \item \textbf{La Sfida Dimensionale:} I problemi del machine learning non sono teoricamente intrattabili: sono quasi sempre formulazioni lineari o convesse. \textbf{La loro difficolt\`a risiede esclusivamente nella scala colossale} (\textit{«hard because large, not hard because hard»}), con milioni o miliardi di variabili.
    \item \textbf{L'Algebra Lineare come Motore dell'Ottimizzazione:} Risolvere problemi su scala titanica richiede l'uso efficiente dell'hardware: l'algebra lineare numerica rappresenta il motore computazionale irrinunciabile che rende l'ottimizzazione rapida, scalabile e priva di instabilit\`a.
\end{enumerate}

\FloatBarrier
\section{Formulario Sinottico dei Modelli Fondamentali}

La tabella~\ref{tab:sinottico-modelli-lezione0} sintetizza i quattro modelli matematici esaminati e le relative propriet\`a computazionali.

\begin{table}[H]
\centering
\footnotesize
\begin{tabularx}{\textwidth}{@{} >{\bfseries}l >{\raggedright\arraybackslash}p{3.8cm} >{\raggedright\arraybackslash}p{3.1cm} Y @{}}
\toprule
\normalfont\textbf{Modello} & \textbf{Formulazione} & \textbf{Paradigma} & \textbf{Note Computazionali} \\
\midrule
Minimi Quadrati & $\min_{\mathbf{w}} \|\mathbf{y} - X\mathbf{w}\|_2^2$ & Equazioni normali:\newline $\mathbf{w}^* = (X^TX)^{-1}X^T\mathbf{y}$ & Risolvere via fattorizzazione QR; mai invertire $X^TX$ esplicitamente. \\
\addlinespace[2pt]
Basso Rango & $\min_{A, B} \|M - AB\|_F^2$ & Decomposizione SVD troncata (Eckart-Young) & Compressione $nm \to k(n+m)$; evitare $M^TM$ per matrici sparse. \\
\addlinespace[2pt]
SVM (Hard) & $\min \frac{1}{2}\|\mathbf{w}_+\|_2^2$\newline $\text{s.t. } y^h(\mathbf{w}_+^T\mathbf{x}^h - w_0) \ge 1$ & Programma quadratico convesso (QP) & Massimizza il margine geometrico $2/\|\mathbf{w}_+\|_2$ a separazione rigida. \\
\addlinespace[2pt]
SVM (Soft) & $\min \frac{1}{2}\|\mathbf{w}_+\|_2^2 + C \sum \xi_h$\newline $\text{s.t. } y^h(\mathbf{w}_+^T\mathbf{x}^h - w_0) \ge 1-\xi_h$ & QP con vincoli lineari (equivalente a Hinge loss) & Parametro $C > 0$ governa il trade-off bias-varianza tollerando violazioni. \\
\addlinespace[2pt]
SVM (Kernel) & $\max_{\boldsymbol{\alpha}} \mathbf{1}^T\boldsymbol{\alpha} - \frac{1}{2}\boldsymbol{\alpha}^T \tilde{K} \boldsymbol{\alpha}$\newline $\text{s.t. } \mathbf{y}^T\boldsymbol{\alpha}=0, \; 0 \le \alpha_h \le C$ & QP con vincoli a scatola (\textit{box constraints}) & Prodotti scalari in spazi di Hilbert $\mathcal{H}$ a dimensione infinita (RBF) in $\BigO{n}$. \\
\addlinespace[2pt]
K-Means & $\min \sum_{i,p} z_{ip} \|\mathbf{c}^p - \mathbf{x}^i\|_2^2$\newline $\text{s.t. } \sum_p z_{ip}=1, \; z_{ip}\in\{0,1\}$ & Minimizzazione alternata (Block Gauss-Seidel) & Convergenza a minimo locale; baricentri calcolati in forma analitica chiusa. \\
\bottomrule
\end{tabularx}
\caption{Formulario sinottico dei modelli matematici introduttivi e delle rispettive propriet\`a computazionali.}
\label{tab:sinottico-modelli-lezione0}
\end{table}
